
% ======================================================================
% Empirical calibration of the effective control affinity
% Requires: amsmath, amssymb, booktabs
% ======================================================================

\subsection{Empirical calibration of the control affinity}
\label{subsec:effective_affinity}

The finite-bias theory describes a controlled interaction through a directional
affinity.  To connect this quantity to the LLM dynamics, we use the microscopic
transitions observed in the state-matched intervention experiment.  We
coarse-grain each focal agent's vote into two states: supporting the controller
target \(Z\), or not supporting \(Z\).  Among controlled updates, we then
estimate
\[
p_{+}
=
P(\text{non-}Z\rightarrow Z\mid \text{controlled}),
\qquad
p_{-}
=
P(Z\rightarrow \text{non-}Z\mid \text{controlled}).
\]
The ratio of these two transition probabilities measures the directional
asymmetry induced by control.  We therefore define the empirical net affinity
as
\[
\boxed{
F_{\mathrm{eff}}
=
\ln\frac{p_{+}}{p_{-}} .
}
\]
A positive \(F_{\mathrm{eff}}\) means that controlled interactions favor
motion toward the controller target, while \(F_{\mathrm{eff}}=0\) corresponds
to equal forward and reverse transition probabilities.

\begin{table}[t]
\centering
\caption{Microscopic transitions observed at controlled update positions.}
\label{tab:effective_affinity_counts}
\begin{tabular}{lrrr}
\toprule
Transition & Eligible updates & Observed transitions & Probability \\
\midrule
non-\(Z\rightarrow Z\) & 572 & 208 & \(p_{+}=0.364\) \\
\(Z\rightarrow\) non-\(Z\) & 508 & 4 & \(p_{-}=0.0079\) \\
\bottomrule
\end{tabular}
\end{table}

Pooling all controlled updates gives
\[
\boxed{
F_{\mathrm{eff}} = 3.83,
}
\]
corresponding to a forward-to-reverse transition ratio
\[
e^{F_{\mathrm{eff}}}\simeq 46.2.
\]
Thus, conditional on observing a controlled microscopic transition
opportunity, motion toward the controller target is approximately \(46\) times
more likely than motion away from it.  A stratified episode-level bootstrap
gives the approximate \(95\%\) interval
\[
\boxed{
F_{\mathrm{eff}}\in[3.06,\,4.88].
}
\]

The transition probabilities also reveal an important kinetic distinction.
In the minimal finite-bias kernel,
\[
p_{+}=\sigma(F),
\qquad
p_{-}=\sigma(-F),
\]
and therefore \(p_{+}+p_{-}=1\).  In the LLM trajectories, however,
\[
p_{+}+p_{-}=0.372.
\]
Most controlled updates therefore leave the focal vote unchanged.  The data
show a strong \emph{directional bias} when revision occurs, together with a
much smaller overall probability of revision.  A convenient way to separate
these two effects is to introduce a kinetic compliance factor \(\gamma\),
\[
p_{+}=\gamma\,\sigma(F),
\qquad
p_{-}=\gamma\,\sigma(-F),
\]
which preserves the same forward-to-reverse affinity while allowing
no-change events.  The empirical value is
\[
\boxed{
\gamma_{\mathrm{eff}}
=
p_{+}+p_{-}
=
0.372.
}
\]

\begin{table}[t]
\centering
\caption{Empirical quantities connecting the microscopic LLM dynamics to the
finite-bias control description.}
\label{tab:effective_affinity_summary}
\begin{tabular}{lll}
\toprule
Quantity & Estimate & Interpretation \\
\midrule
\(F_{\mathrm{eff}}\) & \(3.83\) & net directional affinity \\
\(95\%\) CI for \(F_{\mathrm{eff}}\) & \([3.06,\,4.88]\) & bootstrap uncertainty \\
\(e^{F_{\mathrm{eff}}}\) & \(46.2\) & forward/reverse transition ratio \\
\(\gamma_{\mathrm{eff}}\) & \(0.372\) & probability scale for actual revision \\
\bottomrule
\end{tabular}
\end{table}

These measurements identify the \emph{net} directional affinity of the
controlled LLM dynamics.  They do not by themselves separate the controller
affinity \(h_c\) from the opposing load \(h_e\), since the thermodynamic model
contains only their difference,
\[
F=h_c-h_e,
\]
at the level of controlled transition odds.  Estimating \(h_c\) and \(h_e\)
individually requires an independent characterization of the opposing
epistemic tendency.  We therefore use \(F_{\mathrm{eff}}\) here as the direct
empirical calibration of control strength, while leaving the decomposition
into input and load affinities to the thermodynamic efficiency analysis.
